\def\stat{zat+minin}





\def\tit{ОПРЕДЕЛЕНИЕ ИНДИКАТОРОВ ИНТЕНСИВНОСТИ ПЕРЕНОСА 


НАУЧНЫХ ЗНАНИЙ В СФЕРУ ТЕХНОЛОГИЙ$^*$}





\def\titkol{Определение индикаторов интенсивности переноса 


научных знаний в~сферу технологий}








\def\aut{В.\,А.~Минин$^1$, И.\,М.~Зацман$^2$, В.\,А.~Хавансков$^3$, 


С.\,К.~Шубников$^4$}





\def\autkol{В.\,А.~Минин, И.\,М.~Зацман, В.\,А.~Хавансков, 


С.\,К.~Шубников}





\titel{\tit}{\aut}{\autkol}{\titkol}





\index{Минин В.\,А.}


\index{Зацман И.\,М.}


\index{Хавансков В.\,А.}


\index{Шубников С.\,К.}


\index{Minin V.\,A.}


\index{Zatsman I.\,M.}


\index{Havanskov V.\,A.} 


\index{Shubnikov S.\,K.}








{\renewcommand{\thefootnote}{\fnsymbol{footnote}}


\footnotetext[1]


{Исследование выполнено в~Институте проб\-лем информатики ФИЦ ИУ РАН 


при финансовой поддержке РФФИ 


(проект 16-07-00075).}}





\renewcommand{\thefootnote}{\arabic{footnote}}


\footnotetext[1]{Институт проблем информатики Федерального исследовательского центра 


<<Информатика и~управление>> 


Российской академии наук, \mbox{aleksisss@ya.ru}}


\footnotetext[2]{Институт проблем информатики Федерального исследовательского центра 


<<Информатика и~управление>>


Российской академии наук, \mbox{izatsman@yandex.ru}}


\footnotetext[3]{Институт проблем информатики Федерального исследовательского центра 


<<Информатика и~управление>>


Российской академии наук, \mbox{сhavanskov@yandex.ru}}


\footnotetext[4]{Институт проблем информатики Федерального исследовательского центра 


<<Информатика и~управление>>


Российской академии наук, \mbox{sergeysh50@yandex.ru}}








 \label{st\stat}





                  \thispagestyle{headings}


                  


  \vspace*{-12pt}





     


     





\Abst{Приведены результаты эксперимента по вычислению значений индикаторов 


интенсивности процессов переноса (трансферта) научных знаний в~сферу технологий. 


Эксперимент выполнен с~использованием одиннадцатилетнего массива полнотекстовых 


описаний изобретений. Индикаторное оценивание процесса переноса знаний является 


важным элементом в~исследовании взаимосвязей науки и~технологической сферы, 


характеризующим процессы создания на\-уч\-но-тех\-ни\-че\-ско\-го 


задела в~об\-ласти перспективных 


технологий. В~ходе исследований, проводимых в~ФИЦ <<Информатика и~управ\-ле\-ние>> 


РАН, созданы модель, методика и~макет информационной системы, обеспечивающие 


вычисление значений индикаторов интенсивности переноса  (ИИП) знаний. Макет сис\-те\-мы был 


использован при проведении эксперимента по вычислению индикаторов для сферы 


информационных технологий. Основная цель проведения этого эксперимента заключается 


в~сопоставлении смыслового содержания (семантики) индикаторов двух видов, 


характеризующих разные аспекты интенсивности процессов переноса. Проведено 


сопоставление вы\-чис\-лен\-ных значений этих индикаторов и~уточнено их смысловое 


содержание.}





\KW{информационные взаимосвязи науки и~технологий; цитирование научных работ; 


интенсивность процессов переноса знаний; индикаторное оценивание; информационные 


технологии}





\DOI{10.14357\!/\!08696527180315} 








\vskip 10pt plus 9pt minus 6pt








      \thispagestyle{headings}


      


      \vspace*{-16pt} 


     


\section{Введение}





\vspace*{-4pt}


     


     Цель данной статьи~--- сопоставление и~уточнение смыслового 


содержания индикаторов двух видов, характеризующих разные аспекты 


интенсивности процессов переноса научных знаний в~сферу информационных 


технологий, на основе результатов эксперимента по вычислению их значений. 


Эксперимент проводился в~Институте проблем информатики Федерального 


исследовательского центра <<Информатика и~управление>> РАН в~рамках 


проекта по гранту РФФИ <<Методы и~информационные технологии 


оценивания интенсивности переноса научных знаний в~сферу разработки 


технологий>>.


     


     Целью проекта является разработка методов и~средств информатики для 


исследования этих процессов и~оценки их интенсивности. Это исследование 


является развитием работ Института проблем информатики по мониторингу 


и~оцениванию результативности научной деятельности. Индикаторы 


интенсивности процессов переноса научных знаний в~сферу технологий 


предлагается использовать для оценивания <<результатов влияния>> научной 


деятельности. Эти индикаторы дают возможность определять результативность 


создания на\-уч\-но-тех\-ни\-че\-ско\-го задела для развития сферы 


технологий~[1--8].


     


     Несмотря на многолетние теоретические и~эмпирические исследования, 


которые были нацелены на изучение процессов переноса научных  


знаний~[9--17] задача многоаспектной оценки их интенсивности 


с~использованием индикаторов разных видов является до сих пор актуальной, 


хотя первые работы по этой тематике появились около~30~лет назад. 


Например, Е.~Мэнсфилд, проанализировав данные о взаимосвязи результатов 


научных исследований и~технологических инноваций, пришел к~выводу, что 


<<$\ldots$10\% технологических новшеств не были бы изобретены или же были 


бы созданы с~большой задержкой, если бы они были сделаны без 


использования результатов соответствующих академических 


исследований>>~\cite{16-mm}.


     


     В рамках проекта разработаны модель и~методология формирования 


системы количественных ИИП знаний 


из различных направлений научных исследований (ННИ) в~сферу технологий. 


Модель и~методология основаны на категоризации оснований формирования 


системы индикаторов, включающей~35 их видов~\cite{18-mm, 19-mm}. 


В~частности, к~числу этих оснований, использованных для категоризации 


индикаторов по видам, относятся~\cite{19-mm}: 


     \begin{itemize}


\item характер переносимых знаний (фундаментальные и~прикладные 


исследования); 


     \item способ учета совпадающих рубрик классификаторов ННИ, 


присвоенных научным публикациям, цитируемым в~изобретении (фактический 


учет и~час\-тот\-ный).


     \end{itemize}


     


     Формирование ИИП научных знаний требует автоматизированной 


обработки больших объемов слабоструктурированных полнотекстовых 


описаний изобретений и~сопоставления метаданных патентной информации 


с~метаданными научных публикаций (в журналах или материалах 


конференций), цитируемых в~описаниях изобретений, для чего необходимо 


было разработать аналитическую информационную систему  


(АИС)~\cite{20-mm}.





\vspace*{-4pt}


     


\section{Порядок проведения эксперимента}





\vspace*{-2pt}


     


     В эксперименте для вычисления значений ИИП научных знаний был 


разработан и~использован макет АИС, описанный в~\cite{20-mm}. Макет (см.\ 


рисунок) состоит из двух функциональных подсистем, взаимодействующих 


с~основной базой данных (ОБД):


     \begin{enumerate}[(1)]


\item подсистема лингвистической обработки полнотекстовых описаний 


и~библиографических данных изобретений;


\item подсистема вычисления значений ИИП научных знаний в~сферу 


технологий.


\end{enumerate}





\begin{figure} %fig1


\vspace*{1pt}


 \begin{center}


 \mbox{%


 \epsfxsize=83mm %5.637mm 


 \epsfbox{min-1.eps}


 }





\vspace*{3pt}





{\small Обобщенная структурная схема макета информационной системы}


\end{center}


\vspace*{-18pt}


\end{figure}








     Основная база данных состоит из следующих взаимосвязанных таблиц:


     \begin{enumerate}[(1)]


\item таблица сценариев\footnote{Описание сценариев исследования дано в~работе~\cite{19-mm}.}, содержащая параметры, определяющие 


заданный вариант исследования процессов переноса знаний (ТбС);


\item таблица изобретений, содержащая их библиографические данные; 


каждое изобретение соотнесено с~некоторым 


сценарием\footnote{Возможна привязка одного изобретения к~нескольким 


заданным сценариям, если исследуемые в~них технологические области 


имеют тематическое пересечение.} (ТбПат);


\item таблица ссылок на цитируемые публикации, каждая из которых 


привязана к~определенному изобретению, из полного описания которого 


были извлечены ссылки (ТбСПуб);


\item таблица кодов рубрик ННИ, которые присваиваются цитируемой 


в~изобретении публикации (ТбРуб);








\item таблица кортежей вида $\langle$индекс МПК, ссылка на 


публикацию$\rangle$ (ТбКрПуб);


\item таблица кортежей вида $\langle$индекс МПК, рубрика 


ННИ$\rangle$ (ТбКрРуб).


\end{enumerate}





     Описание упомянутых в~списке таблиц кортежей дано  


в~работе~\cite{19-mm}. 


     


     Кроме вышеперечисленных таблиц ОБД в~состав информационной 


системы входит технологическая база данных (ТБД), которая содержит ряд 


вспомогательных данных:


     \begin{enumerate}[(1)]


\item рубрики классификаторов ННИ, используемых в~макете (ТбКННИ);


\item таблица шаблонов для поиска ссылок на публикации в~текстах 


описаний изобретений (ТбШП);


\item таблица нормализованных названий научных периодических 


изданий с~указанием кодов их тематических рубрик ННИ (ТбНПИ).


     \end{enumerate}


     


     В макете предусмотрено использование двух классификаторов ННИ: 


классификатор Российского фонда фундаментальных исследований (РФФИ) 


и~более прикладной Государственный рубрикатор на\-уч\-но-тех\-ни\-че\-ской 


информации (\mbox{ГРНТИ}).


     


     Эксперимент включает в~себя последовательное выполнение ряда этапов 


обработки первичных и~вторичных информационных массивов. Содержание 


и~объем первичных массивов (включающих библиографические данные 


и~полнотекстовые описания запатентованных изобретений) обусловлен 


исследуемой технологической сферой и~временн$\acute{\mbox{ы}}$м периодом 


исследования трансферта научных знаний в~данную технологическую сферу. 


В~данном эксперименте исследуемой технологической сферой послужили 


информационные технологии, к~которым были отнесены патенты на 


изобретение с~классом МПК (Международной патентной классификации) G06 


<<\textit{Обработка данных}; \textit{вычисление; счет}>> и~двумя 


подклассами: G10L <<\textit{Анализирование или синтезирование речи}; 


\textit{распознавание речи}; \textit{обработка речи или голоса}$\ldots$>> и~G11C 


<<\textit{Запоминающие устройства статического типа}$\ldots$>>. 


Исследуемый в~рамках эксперимента период времени составляет~11~лет: 


с~2002 по~2012~год включительно. В~сценарии этот период может быть задан 


по дате подачи заявок на патентование или публикации патентов, относящихся к~исследуемой технологической сфере.


     


     На первом этапе проведения эксперимента был сформирован массив 


библиографических данных тех патентуемых изобретений, которые содержали 


ссылки на полнотекстовые описания изобретений. Следующим этапом стал 


поиск (в~полнотекстовых описаниях изобретений) ссылок на цитируемые в~них 


публикации, их выделение и~экспертная верификация. В~рамках эксперимента 


рассматривались ссылки на научные публикации в~журналах и~в~трудах 


научных конференций.


     


     Далее для журналов и~трудов конференций, выделенных из массива 


ссылок на научные публикации, экспертами было проведено их (журналов 


и~трудов конференций) рубрицирование по классификаторам ННИ РФФИ 


и~ГРНТИ (использовалась таблица нормализованных названий научных 


периодических изданий). Следует отметить, что в~макете принято следующее 


допущение: научные публикации наследуют рубрики ННИ журналов или 


трудов конференций, в~которых они опубликованы, что обеспечивает 


автоматическое формирование записей в~таблице кодов рубрик ННИ.


     


     Таким образом, по завершении данного этапа в~макете были 


сформированы таблицы двух видов кортежей: $\langle$индекс МПК, ссылка на 


публикацию$\rangle$ и~$\langle$индекс МПК, рубрика ННИ$\rangle$.


     


     Макет АИС включает программу вычисления значений ИИП научных 


знаний. Параметрически для этой программы можно задавать тот уровень 


тематической иерархии МПК и~классификаторов ННИ, для которого 


вычисляются значения ИИП. Наиболее детальная информация об 


интенсивности процессов переноса получается при задании самого нижнего 


уровня иерархии. Если задать более высокий уровень тематической иерархии, 


то получим (из таблицы вида $\langle$индекс МПК, ссылка на 


публикацию$\rangle$ или таблицы вида $\langle$индекс МПК, рубрика 


ННИ$\rangle$) сгруппированные относительно нижних рубрик данные об 


интенсивности процессов переноса знаний. Остановимся более подробно на 


описании уровней тематической иерархии МПК и~классификаторов ННИ 


и~рассмотрим пример группировки данных.


     


     Индекс МПК по своей структуре является позиционным кодом и~состоит 


из последовательности символов, используемых для обозначения раздела, 


класса, подкласса, группы и~подгруппы МПК~\cite{21-mm}. В~качестве 


примера в~табл.~1 описана структура индекса G06F~3\!/\!048 <<Средства 


взаимодействия, основанные на графических интерфейсах пользователя>> 


в~соответствии с~пятью уровнями МПК.


     


     \begin{table}[b]\small %tabl1


     \vspace*{-24pt}


     \begin{center}


     \Caption{Структура индекса МПК G06F 3\!/\!048}


      \vspace*{2ex}


      


      \tabcolsep=3pt


      \begin{tabular}{|c|c|c|c|c|c|}


      \hline


G&06&F&3&/&048\\


\hline


\tabcolsep=0pt\begin{tabular}{c}Раздел~---\\ 1-й уровень\end{tabular}&


\tabcolsep=0pt\begin{tabular}{c}Класс~---\\ 2-й уровень\end{tabular}&


\tabcolsep=0pt\begin{tabular}{c}Подкласс~---\\ 3-й уровень\end{tabular}&


\tabcolsep=0pt\begin{tabular}{c}Группа~---\\ 4-й  уровень\end{tabular}&


\tabcolsep=0pt\begin{tabular}{c}Разделитель\\ между группой\\ и~подгруппой\end{tabular}&


\tabcolsep=0pt\begin{tabular}{c}Подгруппа~---\\  5-й уровень\end{tabular}\\


\hline


\end{tabular}


\end{center}


\end{table}


     


     Коды рубрик ННИ также содержат несколько уровней (как правило, три), 


но их структуры существенно отличаются как от структуры МПК, так и~друг от 


друга. Например, код рубрики классификатора ГРНТИ имеет 6 позиций, при 


этом три уровня рубрикации разделены точками (XX.YY.ZZ). Код рубрики 


классификатора РФФИ включает в~себя три уровня и~5~позиций (XX-YZZ). 


Первый уровень (XX) отделен от второго и~третьего уровней дефисом. Второй 


уровень определяется значением кода рубрики в~первой позиции после 


дефиса~(Y). Третий уровень определяется значениями кода рубрики 


в~последних двух позициях кода (ZZ).


     


     Рассмотрим фрагмент массива кортежей и~их весов\footnote{Определение 


понятия <<вес кортежа>> для классификаторов РФФИ и~ГРНТИ дано 


в~работе~\cite{19-mm}.} (табл.~2) и~результат его обработки, 


иллюстрирующего тематическую группировку данных об интенсивности 


процессов переноса. Рубрики ННИ в~данном примере приведены согласно 


классификатору РФФИ.


     


     \begin{table}\small %tabl2


     \begin{center}


    \parbox{212pt}{ \Caption{Фрагмент массива кортежей вида $\langle$индекс МПК, рубрика 


ННИ$\rangle$}


}





      \vspace*{2ex}


      


      \begin{tabular}{|l|c|c|}


      \hline


\multicolumn{1}{|c|}{Индекс МПК}&Рубрика ННИ&


\tabcolsep=0pt\begin{tabular}{c}Вес кортежа\\ (условно)\end{tabular}\\


\hline


G05B 19\!/\!18&01-215&$V_1$\hphantom{$_1$}\\


G05B 19\!/\!18&07-410&$V_2$\hphantom{$_1$}\\


G05B 19\!/\!18&01-217&$V_3$\hphantom{$_1$}\\


G06F 11\!/\!07&01-215&$V_4$\hphantom{$_1$}\\


G06F 11\!/\!07&07-410&$V_5$\hphantom{$_1$}\\


G06F 11\!/\!07&01-217&$V_6$\hphantom{$_1$}\\


G06F 5\!/\!00&01-215&$V_7$\hphantom{$_1$}\\


G06F 5\!/\!00&07-410&$V_8$\hphantom{$_1$}\\


G06F 5\!/\!00&01-217&$V_9$\hphantom{$_1$}\\


G06K 9\!/\!00&01-214&$V_{10}$\\


G06K 9\!/\!00&01-215&$V_{11}$\\


G06K 9\!/\!00&01-215&$V_{12}$\\


G06K 9\!/\!00&07-410&$V_{13}$\\


G06K 9\!/\!00&01-216&$V_{14}$\\


G06K 9\!/\!00&01-217&$V_{15}$\\


\hline


\end{tabular}


\end{center}


\vspace*{-6pt}


\end{table}





     Если применить к~данному фрагменту массива группировку \textit{на 


уровне подклассов МПК} и~\textit{второго уровня рубрикатора ННИ} (по 


классификатору \mbox{РФФИ}), то получим табл.~3.


     


     \begin{table}\small %tabl3


     \begin{center}


     \parbox{318pt}{\Caption{Пример сгруппированных кортежей $\langle$подкласс МПК\,--\,вто\-рой 


уровень ННИ$\rangle$}


}





      \vspace*{2ex}


      


      \begin{tabular}{|c|c|c|}


      \hline


Индекс МПК&Рубрика ННИ&Суммарные значения весов кортежей\\


\hline


G05B&01-200&$V_1+V_3$\\


G05B&07-400&$V_2$\\


G06F&01-200&$V_4+V_6+V_7+V_9$\\


G06F&07-400&$V_5+V_8$\\


G06K&01-200&$V_{10}+V_{11}+V_{12}+V_{14}+V_{15}$\\


G06K&07-400&$V_{13}$\\


\hline


\end{tabular}


\end{center}


\vspace*{-12pt}


\end{table}


     


     В третьем столбце показаны суммарные значения весов тех кортежей, 


которые соответствуют следующей группировке: подкласс МПК (см.\ 


табл.~1)\,--\,вто\-рой уровень ННИ (по классификатору РФФИ).


     


     Исходные данные о патентах на изобретения и~о~научных публикациях 


изначально (до группировки кортежей) соотнесены с~библиографическими 


данными, относящимися и~к изобретениям, и~к научным публикациям. Поэтому 


кортежи могут быть дополнены, например, некоторыми библиографическими 


данными изобретений и~научных публикаций (табл.~4).


     


     \begin{table}\small %tabl4


     \begin{center}


     \Caption{Кортежи (первые два столбца), дополненные тремя полями 


библиографических данных (последние три столбца)}


      \vspace*{2ex}


      


      \tabcolsep=8pt


      \begin{tabular}{|c|c|c|c|c|}


      \hline


\tabcolsep=0pt\begin{tabular}{c}Индекс МПК\end{tabular}&


Рубрика ННИ&


\tabcolsep=0pt\begin{tabular}{c}Год подачи\\ заявки\end{tabular}&


\tabcolsep=0pt\begin{tabular}{c}Год публикации\\ патента\end{tabular}&


\tabcolsep=0pt\begin{tabular}{c}Год научной\\ публикации\end{tabular}\\


\hline


G05B&01-200&1998&2000&1988\\


G05B&07-400&1998&2000&1991\\


G06F&01-200&1998&2000&1995\\


G06F&07-400&1998&2000&1988\\


G06K&01-200&2000&2001&1986\\


G06K&07-400&2000&2001&1986\\


\hline


\end{tabular}


\end{center}


\vspace*{-12pt}


\end{table}


     


     С помощью SQL-запросов происходит вычисление требуемых значений 


для ИИП научных знаний при заданном уровне группировки. Вычисленные 


значения сохраняются для последующего формирования табличного и\!/\!или 


графического представления значений ИИП научных знаний.





\vspace*{-4pt}





\section{Таблицы значений индикаторов}





\vspace*{-2pt}


     


     Традиции цитирования научных публикаций как в~разных областях 


знаний, так и~в разных сферах технологий заметно различаются. Например, 


библиографические пристатейные списки в~биологии традиционно составляют 


в~среднем не менее~25~ссылок, тогда как, например, в~физике~--- менее~20. 


Поэтому было введено основание категоризации ИИП, описанное 


в~работе~\cite{19-mm}, названное <<Способ учета рубрик классификатора>>, 


которое делит все виды индикаторов на два класса (с~точки зрения способа 


учета рубрик ННИ): фактический и~частотный. Такое деление несет 


следующую смысловую нагрузку. 


     


     Число тех публикаций, которые цитируются в~описании одного 


изобретения и~при этом относятся к~\textit{одной и~той же тематической 


рубрике} используемого классификатора, в~изобретениях, принадлежащих 


к~разным сферам технологий, может существенно различаться. Однако 


наличие в~описании изобретения даже только одной такой ссылки уже 


свидетельствует о возможном случае переноса знаний, относящихся к~этой 


рубрике, в~сферу технологий. Если же цитируемых публикаций, имеющих одну и~ту же рубрику классификатора в~одном описании изобретения, несколько, то 


вторую и~последующие ссылки можно рассматривать как подтверждение 


случаев переноса знаний, относящихся к~этой рубрике ННИ. Для сопоставления 


интенсивности цитирования в~изобретениях, относящихся к~\textit{разным 


сферам} технологий (информационные технологии, биотехнологии и~т.\,д.), 


предлагается использовать фактический учет рубрик ННИ, так как в~этом 


случае происходит нормализация по среднему числу тех публикаций, которые 


цитируются в~описании \textit{одного изобретения} и~при этом относятся 


к~\textit{одной и~той же тематической сфере}. Иначе говоря, фактический 


учет руб\-рик ННИ дает возможность проводить сопоставления интенсивности 


в~разных сферах технологий (без учета влияния среднего числа цитируемых 


публикаций в~каждой сфере).


     


     Таким образом, введенное в~работе~\cite{19-mm} основание 


категоризации индикаторов <<Способ учета рубрик классификатора>> состоит в~учете или неучете подтверждений случаев переноса знаний, относящихся 


к~одной рубрике ННИ, что и~дает возможность определить два класса 


индикаторов по этому основанию: \textit{фактический}, когда учитывается 


только факт наличия той или иной рубрики классификатора ННИ у одной или 


нескольких публикаций без учета числа повторений, а также 


\textit{частотный}, когда учитывается число повторений в~одном изобретении. 


Вычисленные значения индикаторов фактического и~частотного классов, 


характеризующие интенсивность переноса знаний из научных дисциплин, 


обозначенных рубриками классификатора РФФИ, в~сферу информационных 


технологий приведены в~табл.~5.


     


     \begin{table}\footnotesize %tabl5


     \vspace*{-4.02498pt}


          \begin{center}


     \Caption{Значения ИИП для двух способов учета рубрик}


      \vspace*{2ex}


      


     \tabcolsep=4pt


      \begin{tabular}{|c|l|c|c|c|c|}


      \hline


\multicolumn{2}{|c|}{Классификатор РФФИ}&\multicolumn{4}{c|}


{Способ учета рубрик классификатора РФФИ}\\


\hline


\multicolumn{1}{|c|}{\raisebox{-12pt}[0pt][0pt]


{\tabcolsep=0pt\begin{tabular}{c}Код\\ рубрики\end{tabular}}}&


\multicolumn{1}{c|}{\raisebox{-12pt}[0pt][0pt]{Название рубрики}}&\multicolumn{2}{c|}{Фактический учет}&\multicolumn{2}{c|}{Частотный 


учет}\\


\cline{3-6}


&&\tabcolsep=0pt\begin{tabular}{c}Абсолютное\\ значение\\ ИИП\end{tabular}&


\tabcolsep=0pt\begin{tabular}{c}Доля\\ рубрики,\\ \%\end{tabular}&


\tabcolsep=0pt\begin{tabular}{c}Абсолютное\\ значение\\ ИИП\end{tabular}&


\tabcolsep=0pt\begin{tabular}{c}Доля\\ рубрики,\\ \%\end{tabular}\\


\hline


01-200&ИНФОРМАТИКА&146\hphantom{9}&19,05&52,66&23,47\hphantom{9}\\


\hline


07-500&\tabcolsep=0pt\begin{tabular}{l}СИСТЕМЫ ОБРАБОТКИ\\ И~АНАЛИЗА ДАННЫХ


\end{tabular}&97&12,66&30,95&13,79\hphantom{9}\\


\hline


07-800&\tabcolsep=0pt\begin{tabular}{l}ФУНДАМЕНТАЛЬНЫЕ\\ ОСНОВЫ\\


 ИНФОРМАЦИОННЫХ\\ ТЕХНОЛОГИЙ,\\ ВЫЧИСЛИТЕЛЬНЫХ\\ 


СИСТЕМ\\ И ТЕЛЕКОММУНИКАЦИЙ\end{tabular}&89&11,61&19,14&8,54\\


\hline


07-700&\tabcolsep=0pt\begin{tabular}{l}ЭЛЕМЕНТНАЯ БАЗА\\ ИНФОРМАЦИОННЫХ\\ СИСТЕМ\end{tabular}&33&


\hphantom{9}4,30&18,49&8,24\\


\hline


08-600&\tabcolsep=0pt\begin{tabular}{l}ТЕХНИЧЕСКИЕ СИСТЕМЫ\\ И ПРОЦЕССЫ\\ УПРАВЛЕНИЯ\end{tabular}&53&


\hphantom{9}6,91&15,85&7,06\\


\hline


07-600&\tabcolsep=0pt\begin{tabular}{l}ВЫЧИСЛИТЕЛЬНЫЕ\\ И СЕТЕВЫЕ РЕСУРСЫ\end{tabular}&


46&\hphantom{9}6,00&13,02&5,80\\


\hline


07-400&\tabcolsep=0pt\begin{tabular}{l}ИНФОРМАЦИОННЫЕ\\ РЕСУРСЫ И СИСТЕМЫ\end{tabular}&


58&\hphantom{9}7,57&\hphantom{9}9,92&4,42\\


\hline


01-100&МАТЕМАТИКА&50&\hphantom{9}6,52&\hphantom{9}9,60&4,28\\


\hline


02-400&\tabcolsep=0pt\begin{tabular}{l}РАДИОФИЗИКА,\\ ЭЛЕКТРОНИКА,\\ АКУСТИКА


\end{tabular}&35&\hphantom{9}4,56&\hphantom{9}9,48&4,22\\


\hline


02-200&\tabcolsep=0pt\begin{tabular}{l}ФИЗИКА\\ КОНДЕНСИРОВАННЫХ\\ СРЕД\end{tabular}&


34&\hphantom{9}4,43&\hphantom{9}8,87&3,95\\


\hline


08-300&\tabcolsep=0pt\begin{tabular}{l}ЭЛЕКТРОФИЗИКА,\\ ЭЛЕКТРОТЕХНИКА\\ 


И ЭЛЕКТРОЭНЕРГЕТИКА\end{tabular}&13&\hphantom{9}1,69&\hphantom{9}8,68&3,87\\


\hline


02-300&\tabcolsep=0pt\begin{tabular}{l}ОПТИКА, КВАНТОВАЯ\\ ЭЛЕКТРОНИКА\end{tabular}&


27&\hphantom{9}3,52&\hphantom{9}7,05&3,14\\


\hline


04-300&\tabcolsep=0pt\begin{tabular}{l}ФУНДАМЕНТАЛЬНАЯ\\ МЕДИЦИНА\\ И ФИЗИОЛОГИЯ


\end{tabular}&10&\hphantom{9}1,30&\hphantom{9}4,90&2,18\\


\hline


05-700&\tabcolsep=0pt\begin{tabular}{l}ГЕОГРАФИЯ\\ И ГИДРОЛОГИЯ СУШИ\end{tabular}&12&


\hphantom{9}1,56&\hphantom{9}3,47&1,55\\


06-200&\tabcolsep=0pt\begin{tabular}{l}ЭКОНОМИЧЕСКИЕ\\ НАУКИ\end{tabular}&4&


\hphantom{9}0,52&\hphantom{9}2,10&0,94\\


\hline


05-400&ГЕОФИЗИКА&6&\hphantom{9}0,78&\hphantom{9}1,97&0,88\\


\hline


06-300&\tabcolsep=0pt\begin{tabular}{l}ФИЛОСОФСКИЕ\\ И СОЦИАЛЬНЫЕ НАУКИ,\\ ПСИХОЛОГИЯ\end{tabular}&


12\hphantom{9}&\hphantom{9}1,56&\hphantom{9}1,60&0,71\\


\hline


\multicolumn{6}{r}{\textit{Окончание табл.~5  на с.}~\pageref{tttr}}


\end{tabular}


\end{center}


\end{table}











     


     В заключение этого раздела отметим важный терминологический аспект в~названиях рубрик классификаторов. Слово <<информатика>> 


в~классификаторе РФФИ и~рубрикаторе ГРНТИ обозначает две разные 


научные дисциплины. В~классификаторе РФФИ это слово обозначает 


компьютерную науку (computer science), а~в~ГРНТИ~--- информационную 


науку (information science), которая на современном этапе является 


профильной, но не базовой научной дисциплиной для создания 


информационных технологий.





\vspace*{-4pt}


     


\section{Заключение}





\vspace*{-2pt}


     


     Проведенный эксперимент подтвердил, что разработанные в~ходе 


выполнения проекта РФФИ 16-07-00075 и~использованные в~данном 


эксперименте концепция и~архитектура информационной системы, как 


и~проектные решения функциональных подсистем, дают возможность вычислить 


значения ИИП знаний из различных ННИ, что 


позволяет экспертам оценить количественно степень воздействия различных 


ННИ на развитие той или иной сферы технологий.


     


     Введенное в~работе~\cite{19-mm} основание категоризации индикаторов 


<<Способ учета рубрик классификатора>> впервые дало возможность 


определить два класса индикаторов по этому основанию: фактический, когда 


учитывается только факт наличия той или иной рубрики классификатора ННИ 


у одной или нескольких публикаций без учета числа повторений, и~частотный, 


когда учитывается число повторений в~одном изобретении.





\setcounter{table}{4}


\begin{table}\footnotesize %tabl5


     \begin{center}


     \Caption{(\textit{окончание}) Значения ИИП для двух способов учета рубрик}


     \label{tttr}


      \vspace*{2ex}


      


     \tabcolsep=4.5pt


      \begin{tabular}{|c|l|c|c|c|c|}


      \hline


\multicolumn{2}{|c|}{Классификатор РФФИ}&\multicolumn{4}{c|}


{Способ учета рубрик классификатора РФФИ}\\


\hline


\multicolumn{1}{|c|}{\raisebox{-12pt}[0pt][0pt]


{\tabcolsep=0pt\begin{tabular}{c}Код\\ рубрики\end{tabular}}}&


\multicolumn{1}{c|}{\raisebox{-12pt}[0pt][0pt]{Название рубрики}}&\multicolumn{2}{c|}{Фактический учет}&\multicolumn{2}{c|}{Частотный 


учет}\\


\cline{3-6}


&&\tabcolsep=0pt\begin{tabular}{c}Абсолютное\\ значение\\ ИИП\end{tabular}&


\tabcolsep=0pt\begin{tabular}{c}Доля\\ рубрики,\\ \%\end{tabular}&


\tabcolsep=0pt\begin{tabular}{c}Абсолютное\\ значение\\ ИИП\end{tabular}&


\tabcolsep=0pt\begin{tabular}{c}Доля\\ рубрики,\\ \%\end{tabular}\\


\hline


04-200&\tabcolsep=0pt\begin{tabular}{l}ФИЗИКО-ХИМИЧЕСКАЯ\\ БИОЛОГИЯ\end{tabular}&6&0,78&1,52&0,68\\


\hline


02-100&ЯДЕРНАЯ ФИЗИКА&5&0,65&0,83&0,37\\


\hline


02-900&\tabcolsep=0pt\begin{tabular}{l}МЕДИЦИНСКАЯ\\ ФИЗИКА\end{tabular}&9&1,17&0,72&0,32\\


\hline


05-500&ОКЕАНОЛОГИЯ&2&0,26&0,62&0,27\\


\hline


08-100&\tabcolsep=0pt\begin{tabular}{l}МАШИНОВЕДЕНИЕ\\ И ИНЖЕНЕРНАЯ\\ МЕХАНИКА\end{tabular}&3&0,39&0,61&0,27\\


\hline


02-600&ФИЗИКА ПЛАЗМЫ&2&0,26&0,57&0,25\\


\hline


03-200&\tabcolsep=0pt\begin{tabular}{l}НЕОРГАНИЧЕСКАЯ\\ ХИМИЯ\end{tabular} &1&0,13&0,40&0,18\\


\hline


04-100&ОБЩАЯ БИОЛОГИЯ&6&0,78&0,40&0,18\\


\hline


06-400&\tabcolsep=0pt\begin{tabular}{l}ФИЛОЛОГИЯ,\\   КУЛЬТУРОЛОГИЯ,\\


 ИСКУССТВОЗНАНИЕ,\\   ОХРАНА ПАМЯТНИКОВ\\ ИСТОРИИ И 


КУЛЬТУРЫ\end{tabular}&2&0,26&0,33&0,15\\


\hline


05-600&ФИЗИКА АТМОСФЕРЫ&2&0,26&0,31&0,14\\


\hline


02-800&АСТРОНОМИЯ&3&0,39&0,30&0,13\\


\hline


02-700&\tabcolsep=0pt\begin{tabular}{l}ТЕОРЕТИЧЕСКАЯ\\ ФИЗИКА\end{tabular}&1&0,13&0,06&0,02\\


\hline


\end{tabular}


\end{center}


\vspace*{-9pt}


\end{table}


     


     Фактический способ учета позволяет:


     \begin{itemize}


     \item провести нормализацию по среднему числу тех публикаций, 


которые цитируются в~описании одного изобретения и~при этом относятся 


к~одной и~той же тематической рубрике ННИ;


     \item сопоставить интенсивности цитирования в~изобретениях, 


относящихся к~разным сферам технологий, без учета влияния среднего числа 


ссылок.


     \end{itemize}


     


     Частотный способ учета позволяет сопоставлять интенсивности 


цитирования в~изобретениях с~учетом влияния среднего числа ссылок. Такой 


учет, согласно табл.~5, для одних ННИ приводит к~уменьшению относительных 


значений ИИП частотного класса по сравнению с~ИИП фактического класса, 


а~других ННИ~--- к~их увеличению (см., например, рубрику~08-300, для 


которой частотный способ учета увеличивает долю этой рубрики более чем 


в~2~раза).





\vspace*{-4pt}





{\small\frenchspacing


{%\baselineskip=10.8pt


\addcontentsline{toc}{section}{Литература}


\begin{thebibliography}{99}





\vspace*{-2pt}





\bibitem{1-mm}


\Au{Зацман И.\,М., Веревкин Г.\,Ф.} Информационный мониторинг сферы науки в~задачах 


программно-целевого управления~/\!\!/ Системы и~средства информатики, 2006. Т.~16. №\,1. 


С.~185--210.


\bibitem{2-mm}


\Au{Зацман И.\,М., Шубников С.\,К.} Принципы обработки информационных ресурсов для 


оценки инновационного потенциала направлений научных исследований~/\!\!/ Электронные 


библиотеки: перспективные методы и~технологии, электронные коллекции: Тр. IX~Всеросс. 


научн. конф. RCDL'2007.~--- Переславль: Университет города Переславля, 2007. С.~35--44.


\bibitem{3-mm}


\Au{Зацман И.\,М., Кожунова О.\,С.} Семантический словарь системы информационного 


мониторинга в~сфере науки: задачи и~функции~/\!\!/ Системы и~средства информатики, 2007. 


Т.~17. №\,1. С.~124--141.





\bibitem{7-mm} %4


\Au{Zatsman I., Kozhunova O.} Evaluating for institutional academic activities: Classification 


scheme for R\&D indicators~/\!\!/ 10th Conference (International) 


on Science and Technology 


Indicators: Book of Abstracts.~--- Vienna: ARC GmbH, 2008. P.~428--431.


\bibitem{8-mm} %5


\Au{Zatsman I., Kozhunova O.} Evaluation system for the Russian Academy of Sciences: 


Objectives--resources--results approach and R\&D indicators~/\!\!/  Atlanta Conference on 


Science and Innovation Policy Proceedings~/ Eds. S.\,E.~Cozzens, P.~Catalаn.~---


Atlanta, GA, USA, 2009.


{\sf 


http://smartech.gatech.edu/bitstream/1853/32300/1/104-674-1-PB.pdf}.


\bibitem{4-mm} %6


\Au{Zatsman I., Durnovo A.} Incompleteness problem of indicators system of research 


programme~/\!\!/ 11th Conference (International) on Science and Technology Indicators: 


Book of Abstracts.~--- Leiden: CWTS, 2010. P.~309--311.


\bibitem{5-mm} %7


\Au{Архипова М.\,Ю., Зацман И.\,М., Шульга~С.\,Ю.} Индикаторы патентной активности 


в~сфере информационно-коммуникационных технологий и~методика их вычисления~/\!\!/ 


Экономика, статистика и~информатика. Вестник УМО, 2010. №\,4. С.~93--104.


\bibitem{6-mm} %8


\Au{Зацман И.\,М., Дурново А.\,А.} Моделирование процессов формирования экспертных 


знаний для мониторинга про\-грам\-мно-це\-ле\-вой деятельности~/\!\!/ Информатика и~её 


применения, 2011. Т.~5. Вып.~4. С.~84--98.











\bibitem{12-mm} %9


\Au{Narin F., Noma E.} Is technology becoming science?~/\!\!/ Scientometrics, 1985. Vol.~7. 


No.\,3--6. P.~369--381.





\bibitem{14-mm} %10


\Au{Mansfield E.} Academic research and innovation~/\!\!/ Res. Policy, 1991. Vol.~20. No.\,1. 


P.~1--12.


\bibitem{9-mm} %11


\Au{Schmoch U.} Tracing the knowledge transfer from science to technology as reflected in patent 


indicators~/\!\!/ Scientometrics, 1993. Vol.~26. P.~193--211.


\bibitem{15-mm} %12


\Au{Mansfield E.} Academic research underlying industrial innovations: Sources, characteristics 


and financing~/\!\!/ Rev. Econ. Statistics, 1995. Vol.~77. No.\,1. P.~55--62.


\bibitem{13-mm} %13


\Au{Narin F., Olivastro D.} Linkage between patents and papers: An interim EPO/US 


comparison~/\!\!/ Scientometrics, 1998. Vol.~41. No.\,1-2. P.~51--59.


\bibitem{16-mm} %14


\Au{Mansfield E.} Academic research and industrial innovation: An update of empirical 


findings~/\!\!/ Res. Policy, 1998. Vol.~26. No.\,7-8. P.~773--776.





\bibitem{10-mm} %15


\Au{Tijssen R.\,J.\,W., Buter R.\,K., Van Leeuwen~Th.\,N.} Technological relevance of science: An 


assessment of citation linkages between patents and research papers~/\!\!/ Scientometrics, 2000. 


Vol.~47. No.\,2. P.~389--412.


\bibitem{11-mm} %16


\Au{Van Looy B., Zimmermann E., Veugelers~R., Verbeek~A., Mello~J., Debackere~K.} Do 


science--technology interactions pay on when developing technology? An exploratory investigation 


of~10 science-intensive technology domains~/\!\!/ Scientometrics, 2003. Vol.~57. No.\,3.  


P.~355--367.


\bibitem{17-mm}


Third European Report on Science \& Technology Indicators.~--- Luxembourg: Office for Official 


Publications of the European Communities, 2003. 451~p.


\bibitem{18-mm}


\Au{Минин В.\,А., Зацман И.\,М., Кружков~М.\,Г., Норекян~Т.\,П.} Методологические 


основы создания информационных систем для вычисления индикаторов тематических 


взаимосвязей науки и~технологий~/\!\!/ Информатика и~её применения, 2013. Т.~7. Вып.~1. 


С.~70--81.


\bibitem{19-mm}


\Au{Минин В.\,А., Зацман И.\,М., Хавансков~В.\,А., Шубников~С.\,К.} Методы индикаторного 


оценивания процессов переноса знаний из области научных исследований в~сферу 


технологического развития~/\!\!/ На\-уч\-но-тех\-ни\-че\-ская информация. Сер.~2, 2017. №~8. 


С.~1--12.


\bibitem{20-mm}


\Au{Минин В.\,А., Зацман~И.\,М., Хавансков~В.\,А., Шубников~С.\,К.} Макетирование 


информационной системы вычисления индикаторов переноса научных знаний в~сферу 


технологий~/\!\!/ Система и~средства информатики, 2017. Т.~27. №\,3. С.~171--182.


\bibitem{21-mm}


Архитектоника классификационных индексов. 


{\sf http://www1.fips.ru/wps/wcm/\linebreak 


connect/content\_ru/ru/inform\_resources/international\_classification/inventions/\linebreak


mpk\_begin/article\_2.}


        \end{thebibliography}


} }








\vspace*{-3pt}





\hfill{\small\textit{Поступила в~редакцию 10.08.18}}














\vspace*{8pt}





\hrule





\vspace*{2pt}





%\newpage





%\vspace*{-28pt}





\hrule





%\vspace*{9pt}











\def\tit{CALCULATION OF INTENSITY INDICATORS  


OF~SCIENCE-TECHNOLOGY LINKAGES}





\def\titkol{Calculation of intensity indicators  


of~science-technology linkages}








\def\aut{V.\,A.~Minin, I.\,M.~Zatsman, V.\,A.~Havanskov, 


and~S.\,K.~Shubnikov}


\def\autkol{V.\,A.~Minin, I.\,M.~Zatsman, V.\,A.~Havanskov, 


and~S.\,K.~Shubnikov}











\titel{\tit}{\aut}{\autkol}{\titkol}





\vspace*{-9pt}





\noindent


      Institute of Informatics Problems, Federal Research Center ``Computer 


Science and~Control'' of the Russian Academy of Sciences, 44-2~Vavilov Str., 


Moscow 119333, Russian Federation








 \noindent





\def\leftfootline{\small{\textbf{\thepage}


\hfill Sistemy i Sredstva Informatiki~---


Systems and Means of Informatics 2018 vol~28 no\,3}


}%


 \def\rightfootline{\small{Sistemy i Sredstva Informatiki~---


 Systems and Means of Informatics   2018   vol~28   no\,3


\hfill \textbf{\thepage}}}


      











\Abste{The outcomes of the experiment on calculating the values of the intensity indicators of 


science-technology linkages are presented. The experiment was carried out using an eleven-year 


array of full-text descriptions of inventions. The indicator evaluation of the process of knowledge 


transfer to the sphere of technology is an important element in the study of the linkages between 


science and technology, characterizing the processes of creating scientific foundations for 


developing advanced technologies. In the course of research conducted at the Institute of 


Informatics Problems of the Russian Academy of Sciences, a~model, a~methodology, and a~testbed of 


an information system have been created for calculation of the values of the intensity indicators of 


science-technology linkages. The main aim of this experiment is to compare the meanings of 


indicators of two types that characterize different aspects of the intensity of knowledge transfer to 


the sphere of technology. Comparison of the calculated values of these indicators is made and their 


meanings are specified.}


     


\KWE{science-technology linkages; citation of publications in inventions; intensity of knowledge 


transfer process; indicator assessment; information technologies}


     


    


      


      


\DOI{10.14357\!/\!08696527180315} 





%\vspace*{-24pt}





\Ack


     \noindent


     The research was supported by the Russian Foundation for Basic Research 


(project 16-07-00075).


 








\renewcommand{\bibname}{\protect\rmfamily References}


%\renewcommand{\bibname}{\large\protect\rm References}





{\small\frenchspacing


{%\baselineskip=10.8pt


\addcontentsline{toc}{section}{References}


\begin{thebibliography}{99}


\bibitem{1-mm-1}


\Aue{Zatsman, I.\,M., and G.\,F.~Verevkin.} 2006. Informatsionnyy monitoring sfery nauki 


v~zadachakh programmno-tselevogo upravleniya [Information monitoring of scientific activities in 


projections tasks with specific objectives]. \textit{Sistemy i~Sredstva Informatiki~--- Systems and 


Means of Informatics} 16(1):185--210.


\bibitem{2-mm-1}


\Aue{Zatsman, I.\,M., and S.\,K.~Shubnikov.} 2007. Printsipy obrabotki informatsionnykh 


resursov dlya otsenki innovatsionnogo potentsiala napravleniy nauchnykh issledovaniy [The 


principles of processing of information resources for estimation of innovative potential of the 


directions of scientific research]. \textit{Tr. IX Vseross. nauchn. konf. RCDL'2007 ``Elektronnye 


biblioteki: perspektivnye metody i~tekhnologii, elektronnye kollektsii''} [9th All-Russian Scientific 


Conference RCDL'2007 ``Digital Libraries: Advanced Methods and Technologies, Digital 


Collections'' Proceedings]. Pereslavl-Zalessky: Pereslavl University. 35--44.


\bibitem{3-mm-1}


\Aue{Zatsman, I.\,M., and O.\,S.~Kozhunova.} 2007. Semanticheskiy slovar' sistemy 


informatsionnogo monitoringa v sfere nauki: zadachi i~funktsii [Semantic dictionary of information 


monitoring system in science: Tasks and functions]. \textit{Sistemy i~Sredstva Informatiki~--- 


Systems and Means of Informatics} 17(1):124--141.


\bibitem{7-mm-1} %4


\Aue{Zatsman, I., and O. Kozhunova.} 2008. Evaluating for institutional academic activities: 


Classification scheme for R\&D indicators. \textit{10th  Conference (International) on Science and 


Technology Indicators: Book of Abstracts}. Vienna: ARC GmbH. 428--431.


\bibitem{8-mm-1} %5


\Aue{Zatsman, I., and O.~Kozhunova.} 2009.


Evaluation system for the Russian Academy of Sciences: 


Objectives--resources--results approach and R\&D indicators. \textit{Atlanta Conference on 


Science and Innovation Policy Proceedings}. Eds. S.\,E.~Cozzens and P.~Catalаn. 


Atlanta, GA.


Available at: {\sf 


http://smartech.gatech.edu/bitstream/1853/\linebreak 32300/1/104-674-1-PB.pdf} (accessed August~10, 2018).





\bibitem{4-mm-1} %6


\Aue{Zatsman, I., and A. Durnovo.} 2010. Incompleteness problem of indicators system of 


research programme. \textit{11th Conference (International) 


on Science and Technology Indicators: Book of Abstracts}. Leiden: CWTS. 309--311.


\bibitem{5-mm-1} %7


\Aue{Arkhipova, M.\,Yu., I.\,M. Zatsman, and S.\,Yu.~Shul'ga.} 2010. Indikatory patentnoy 


aktivnosti v~sfere informatsionno-kommunikatsionnykh tekhnologiy i~metodika ikh vychisleniya 


[Indicators of patent activities in the sphere of information and communication technologies and a 


technique of their computation]. \textit{Ekonomika, statistika i~informatika. Vestnik UMO} 


[Economics, Statistics, and Informatics. UMO Bull.] 4:93--104.


\bibitem{6-mm-1} %8


\Aue{Zatsman, I.\,M., and A.\,A.~Durnovo.} 2011. Modelirovanie protsessov formirovaniya 


ekspertnykh znaniy dlya monitoringa programmno-tselevoy deyatel'nosti [Modeling of processes 


for creation of expert knowledge for monitoring of goal-oriented programme activities]. 


\textit{Informatika i~ee Primeneniya~--- Inform. Appl.} 5(4):84--98.








\bibitem{12-mm-1} %9


\Aue{Narin, F., and E. Noma.} 1985. Is technology becoming science? \textit{Scientometrics} 


7(3--6):369--381.


\bibitem{14-mm-1} %10


\Aue{Mansfield, E.} 1991. Academic research and innovation. \textit{Res. Policy}  


20(1):1--12.


\bibitem{9-mm-1} %11


\Aue{Schmoch, U.} 1993. Tracing the knowledge transfer from science to technology as reflected 


in patent indicators. \textit{Scientometrics} 26:193--211.


\bibitem{15-mm-1} %12


\Aue{Mansfield, E.} 1995. Academic research underlying industrial innovations: Sources, 


characteristics and financing. \textit{Rev. Econ. Statistics} 77(1):55--62.


\bibitem{13-mm-1} %13


\Aue{Narin, F., and D. Olivastro.} 1998. Linkage between patents and papers: An


 interim EPO\!/\!US 


comparison. \textit{Scientometrics} 41(1-2):51--59.


\bibitem{16-mm-1} %14


\Aue{Mansfield, E.} 1998. Academic research and industrial innovation: An update of empirical 


findings. \textit{Res. Policy} 26(7-8):773--776.


\bibitem{10-mm-1} %15


\Aue{Tijssen, R.\,J.\,W., R.\,K.~Buter, and Th.\,N.~Van Leeuwen.} 2000. Technological relevance 


of science: An assessment of citation linkages between patents and research papers. 


\textit{Scientometrics} 47(2):389--412.


\bibitem{11-mm-1} %16


\Aue{Van Looy, B., E.~Zimmermann, R.~Veugelers, A.~Verbeek, J.~Mello, and K.~Debackere.} 


2003. Do science--technology interactions pay on when developing technology? An exploratory 


investigation of 10 science-intensive technology domains. \textit{Scientometrics}  


57(3):355--367.


\bibitem{17-mm-1}


European Communities. 2003. Third European Report on Science \& Technology Indicators. 


Luxembourg: Office for Official Publications of the European Communities. 451~p.


\bibitem{18-mm-1}


\Aue{Minin, V.\,A., I.\,M. Zatsman, M.\,G.~Kruzhkov, and T.\,P.~Norekyan.} 2013. 


Me\-to\-do\-lo\-gi\-che\-skie osnovy sozdaniya informatsionnykh sistem dlya vychisleniya indikatorov 


tematicheskikh vzaimosvyazey nauki i tekhnologiy [Methodological bases for creating information 


systems calculating indicators of thematic linkages between science and technologies]. 


\textit{Informatika i~ee Primeneniya~--- Inform. Appl.} 7(1):70--81.


\bibitem{19-mm-1}


\Aue{Minin, V.\,A., I.\,M. Zatsman, V.\,A.~Khavanskov, and S.\,K.~Shubnikov.} 2017. 


Methods of indicator-based assessment of knowledge transfer from 


science to technology. \textit{Automatic Documentation Math.


Linguistics}  %[Scientific and Technical Information. Series~2] 


51(4):180--190. %(8):1--12.


\bibitem{20-mm-1}


\Aue{Minin, V.\,A., I.\,M. Zatsman, V.\,A.~Havanskov, and S.\,K.~Shubnikov.} 2017. 


Maketirovaniye informatsionnoy sistemy vychisleniya indikatorov perenosa nauchnykh znaniy 


v~sferu tekhnologiy [Prototyping an information system for the calculation of indicators of the 


intensity of the transfer of scientific knowledge in the field of technology]. \textit{Sistemy 


i~Sredstva Informatiki~--- Systems and Means of Informatics} 27(3):171--182.


\bibitem{21-mm-1}


Arkhitektonika klassifikatsionnykh indeksov [Architectonics of classification indices]. 


Available at: {\sf 


http://www1.fips.ru/wps/wcm/connect/content\_ru/ru/inform\_\linebreak resources/international\_classification/inventions/mpk\_begin/article\_2} (accessed 


August~10, 2018).


    \end{thebibliography}


} }








\vspace*{-3pt}





\hfill{\small\textit{Received August 10, 2018}}      


      


      \Contr


      


\noindent


\textbf{Minin Vladimir A.} (b.\ 1941)~--- Doctor of Science in physics and


mathematics, consultant, Institute of Informatics 


Problems,


Federal Research Center ``Computer Science and Control'' of the Russian Academy of


Sciences, 44-2~Vavilova Str., Moscow 119333, Russian Federation; \mbox{aleksisss@ya.ru} 





\vspace*{3pt}





\noindent


\textbf{Zatsman Igor M.} (b.\ 1952)~--- Doctor of Science in technology, 


Head of Department, Institute of 


Informatics Problems,


Federal Research Center ``Computer Science and Control'' of the Russian Academy of


Sciences, 44-2~Vavilova Str., Moscow 119333, Russian Federation; \mbox{izatsman@yandex.ru}





\vspace*{3pt}





\noindent


\textbf{Havanskov Valerij A.} (b.\ 1950)~--- scientist, Institute of Informatics Problems,


Federal Research Center ``Computer Science and Control'' of the Russian Academy of


Sciences, 44-2~Vavilova Str., Moscow 119333, Russian Federation; \mbox{сhavanskov@yandex.ru}





\vspace*{3pt}





\noindent


\textbf{Shubnikov Sergej K.} (b.\ 1955)~--- senior scientist, Institute of Informatics Problems,


Federal Research Center ``Computer Science and Control'' of the Russian Academy of


Sciences, 44-2~Vavilova Str., Moscow 119333, Russian Federation; \mbox{sergeysh50@yandex.ru}


      


     


     


\label{end\stat}








\renewcommand{\bibname}{\protect\rm Литература}


     


